\chapter{非局部紧空间上的积分}

尽管拓扑学与测度论之间联系的研究可以追溯到实变量函数现代理论的开端，非局部紧拓扑空间上的积分却只是到了最近才以一般的方式建立起来。在叙述先于现有综合的那些工作之前，我们将回顾有关拓扑与测度之间关系的观念演变中的几个步骤。

对勒贝格来说，问题只在于对一个或多个实变量的函数求积分。1913 年，拉东在 \(\mathbf{R}^n\) 上定义了一般测度及相应的积分；这一理论在 Ch. de la Vallée Poussin 的著作{[}82{]}中有详细的阐述，并且处处依赖于欧氏空间的拓扑性质。稍后，1915 年，弗雷歇在{[}115 c{]}中在一个配有 \(\sigma\) 代数的集合上定义了"抽象"测度及关于这些测度的积分；他指出，这样便有可能不使用任何拓扑手段而建立勒贝格理论的主要结果。他援引{[}115 d{]}引言第一部分中的如下话语来论证他的工作："例如，在一个具有无穷多坐标的空间中，分析的种种应用已经导出收敛序列的多种互不等价的定义，当其中一种定义被另一种替代时，这些空间中序列族与可加集合函数的性质不会有任何改变。"弗雷歇的研究由卡拉泰奥多里加以完成，集合函数扩张为测度这一重要定理正是后者的贡献。Saks 的书的开头{[}268{]}对这一观点给出了浓缩的阐述。

局部紧群上哈尔测度的发现（第~231 页及以后）及其随即得到的众多应用，随后韦伊与盖尔范德在调和分析方面的工作，导致 1940 年前后这一观点的深刻改变：在这类问题中，最方便的是把测度视为连续函数空间上的一个线性型。这一方法被迫局限于紧空间或局部紧空间，但对几乎全部应用而言并不构成障碍；更好的是，J. 泰特与 A. 韦伊把调和分析引入 p 进群和阿代尔群，使解析数论得到了引人注目的更新。

扩大这一观点、考虑非局部紧拓扑空间上测度的必要性来自完全另一个方向：概率论一步一步地引向这类空间的研究，并提供了大量非平凡的例子。这些发展对测度论的影响之所以姗姗来迟，其原因也许要在概率论的相对孤立中去寻找——直到不久之前，它一直处于传统数学学科的边缘。

